\textbf{Phần 1. Trắc nghiệm lựa chọn}
\setcounterpageref{ex}{0}

\begin{ex} %[5D7V38]
Một mảnh đất hình chữ nhật có chu vi $84\text{ m}$. Chiều rộng bằng $\dfrac{3}{4}$ chiều dài. Diện tích của mảnh đất đó là
\choice
{$384\text{ m}^2$}
{\True $432\text{ m}^2$}
{$480\text{ m}^2$}
{$504\text{ m}^2$}
\loigiai{
Nửa chu vi mảnh đất là:
$84 : 2 = 42\text{ (m)}$.\\
Chiều rộng mảnh đất là:
$42 : (3 + 4) \times 3 = 18\text{ (m)}$.\\
Chiều dài mảnh đất là:
$42 - 18 = 24\text{ (m)}$.\\
Diện tích mảnh đất là:
$24 \times 18 = 432\text{ (m}^2\text{)}$.\\
\textbf{Đáp án của câu hỏi là B}
}
\end{ex}

\begin{ex} %[5D7V38]
Hai bạn An và Bình có tất cả $45$ viên bi. Nếu An cho Bình $5$ viên bi thì số bi của An vẫn gấp đôi số bi của Bình. Hỏi lúc đầu An có bao nhiêu viên bi?
\choice
{$30\text{ viên}$}
{\True $35\text{ viên}$}
{$40\text{ viên}$}
{$25\text{ viên}$}
\loigiai{
Sau khi An cho Bình $5$ viên bi thì tổng số bi của hai bạn không đổi và bằng $45$ viên.\\
Khi đó số bi của An là:
$45 : (2 + 1) \times 2 = 30\text{ (viên)}$.\\
Lúc đầu An có số viên bi là:
$30 + 5 = 35\text{ (viên)}$.\\
\textbf{Đáp án của câu hỏi là B}
}
\end{ex}

\begin{ex} %[5D7V38]
Tổng của hai số là $126$. Nếu bớt ở số thứ nhất đi $6$ đơn vị thì số thứ nhất bằng $\dfrac{2}{3}$ số thứ hai. Số thứ nhất ban đầu là
\choice
{$48$}
{\True $54$}
{$72$}
{$78$}
\loigiai{
Tổng mới của hai số sau khi bớt là:
$126 - 6 = 120$.\\
Số thứ nhất sau khi bớt là:
$120 : (2 + 3) \times 2 = 48$.\\
Số thứ nhất ban đầu là:
$48 + 6 = 54$.\\
\textbf{Đáp án của câu hỏi là B}
}
\end{ex}

\begin{ex} %[5D7V38]
Hai lớp 5A và 5B quyên góp được tất cả $280$ quyển sách. Biết rằng $\dfrac{1}{3}$ số sách của lớp 5A bằng $\dfrac{1}{4}$ số sách của lớp 5B. Số sách lớp 5B quyên góp được là
\choice
{$120\text{ quyển}$}
{$140\text{ quyển}$}
{\True $160\text{ quyển}$}
{$180\text{ quyển}$}
\loigiai{
Vì $\dfrac{1}{3}$ số sách 5A bằng $\dfrac{1}{4}$ số sách 5B nên số sách 5A gồm 3 phần, số sách 5B gồm 4 phần.\\
Tổng số phần bằng nhau là:
$3 + 4 = 7\text{ (phần)}$.\\
Số sách lớp 5B quyên góp được là:
$280 : 7 \times 4 = 160\text{ (quyển)}$.\\
\textbf{Đáp án của câu hỏi là C}
}
\end{ex}

\begin{ex} %[5D7V38]
Hiện nay tổng số tuổi của mẹ và con là $48$ tuổi. Biết rằng $4$ năm trước tuổi mẹ gấp $4$ lần tuổi con. Tuổi mẹ hiện nay là
\choice
{$32\text{ tuổi}$}
{\True $36\text{ tuổi}$}
{$40\text{ tuổi}$}
{$38\text{ tuổi}$}
\loigiai{
Bốn năm trước, tổng số tuổi của mẹ và con là:
$48 - 4 \times 2 = 40\text{ (tuổi)}$.\\
Tuổi con $4$ năm trước là:
$40 : (4 + 1) = 8\text{ (tuổi)}$.\\
Tuổi con hiện nay là:
$8 + 4 = 12\text{ (tuổi)}$.\\
Tuổi mẹ hiện nay là:
$48 - 12 = 36\text{ (tuổi)}$.\\
\textbf{Đáp án của câu hỏi là B}
}
\end{ex}

\begin{ex} %[5D7V38]
Một kho có $350\text{ kg}$ gạo gồm gạo tẻ và gạo nếp. Biết rằng số gạo nếp bằng $40\%$ số gạo tẻ. Số ki-lô-gam gạo tẻ là
\choice
{$100\text{ kg}$}
{$150\text{ kg}$}
{\True $250\text{ kg}$}
{$200\text{ kg}$}
\loigiai{
Đổi: $40\% = \dfrac{2}{5}$.\\
Tổng số phần bằng nhau là:
$2 + 5 = 7\text{ (phần)}$.\\
Số gạo tẻ là:
$350 : 7 \times 5 = 250\text{ (kg)}$.\\
\textbf{Đáp án của câu hỏi là C}
}
\end{ex}

\begin{ex} %[5D7V38]
Một sợi dây dài $72\text{ cm}$ được cắt thành hai đoạn, trong đó đoạn thứ nhất dài gấp 3 lần đoạn thứ hai. Chiều dài đoạn thứ nhất là
\choice
{$18\text{ cm}$}
{$36\text{ cm}$}
{\True $54\text{ cm}$}
{$48\text{ cm}$}
\loigiai{
Tổng số phần bằng nhau là:
$3 + 1 = 4\text{ (phần)}$.\\
Chiều dài đoạn thứ nhất là:
$72 : 4 \times 3 = 54\text{ (cm)}$.\\
\textbf{Đáp án của câu hỏi là C}
}
\end{ex}

\begin{ex} %[5D7V38]
Hai số tự nhiên có tổng bằng $242$. Biết rằng nếu viết thêm chữ số $0$ vào bên phải số bé thì được số lớn. Số lớn là
\choice
{$22$}
{\True $220$}
{$200$}
{$210$}
\loigiai{
Viết thêm chữ số $0$ vào bên phải số bé thì được số lớn nên số lớn gấp 10 lần số bé.\\
Tổng số phần bằng nhau: $10 + 1 = 11\text{ (phần)}$.\\
Số lớn là:
$242 : 11 \times 10 = 220$.\\
\textbf{Đáp án của câu hỏi là B}
}
\end{ex}

\begin{ex} %[5D7V38]
Chu vi một hình chữ nhật là $96\text{ cm}$. Nếu tăng chiều rộng thêm $4\text{ cm}$ và giữ nguyên chiều dài thì hình chữ nhật trở thành hình vuông. Chiều rộng hình chữ nhật ban đầu là
\choice
{$20\text{ cm}$}
{\True $22\text{ cm}$}
{$24\text{ cm}$}
{$26\text{ cm}$}
\loigiai{
Nửa chu vi hình chữ nhật là:
$96 : 2 = 48\text{ (cm)}$.\\
Chiều dài hơn chiều rộng là $4\text{ cm}$.\\
Chiều rộng hình chữ nhật là:
$(48 - 4) : 2 = 22\text{ (cm)}$.\\
\textbf{Đáp án của câu hỏi là B}
}
\end{ex}

\begin{ex} %[5D7V38]
Ba tổ công nhân sản xuất được tất cả $180$ sản phẩm. Biết số sản phẩm tổ 1 làm được bằng $\dfrac{1}{2}$ tổng số sản phẩm của tổ 2 và tổ 3. Số sản phẩm tổ 1 làm được là
\choice
{$45\text{ sản phẩm}$}
{\True $60\text{ sản phẩm}$}
{$90\text{ sản phẩm}$}
{$120\text{ sản phẩm}$}
\loigiai{
Coi sản phẩm tổ 1 là 1 phần thì tổng sản phẩm tổ 2 và tổ 3 là 2 phần.\\
Tổng số phần là: $1 + 2 = 3\text{ (phần)}$.\\
Số sản phẩm tổ 1 làm được là:
$180 : 3 \times 1 = 60\text{ (sản phẩm)}$.\\
\textbf{Đáp án của câu hỏi là B}
}
\end{ex}

\begin{ex} %[5D7V38]
Hai số thập phân có tổng bằng $45{,}6$. Số thứ nhất bằng $\dfrac{1}{3}$ số thứ hai. Số thứ hai là
\choice
{$11{,}4$}
{$22{,}8$}
{\True $34{,}2$}
{$30{,}4$}
\loigiai{
Tổng số phần bằng nhau là:
$1 + 3 = 4\text{ (phần)}$.\\
Số thứ hai là:
$45{,}6 : 4 \times 3 = 34{,}2$.\\
\textbf{Đáp án của câu hỏi là C}
}
\end{ex}

\begin{ex} %[5D7V38]
Một trang trại nuôi $450$ con bò gồm bò sữa và bò thịt. Sau khi bán đi $50$ con bò thịt thì số bò thịt còn lại bằng $\dfrac{1}{3}$ số bò sữa. Số bò sữa của trang trại là
\choice
{$100\text{ con}$}
{$200\text{ con}$}
{\True $300\text{ con}$}
{$350\text{ con}$}
\loigiai{
Tổng số bò còn lại là:
$450 - 50 = 400\text{ (con)}$.\\
Tổng số phần bằng nhau là:
$1 + 3 = 4\text{ (phần)}$.\\
Số bò sữa là:
$400 : 4 \times 3 = 300\text{ (con)}$.\\
\textbf{Đáp án của câu hỏi là C}
}
\end{ex}

\textbf{Phần 2. Trắc nghiệm đúng sai}
\setcounterpageref{ex}{0}

\begin{ex} %[5D7V38]
Hai kho thóc chứa tất cả $150$ tấn thóc. Biết rằng số thóc kho thứ nhất bằng $\dfrac{2}{3}$ số thóc kho thứ hai.
\choiceTF
{\True Tổng số phần bằng nhau là $5$ phần}
{\True Kho thứ nhất chứa $60$ tấn thóc}
{\True Kho thứ hai chứa $90$ tấn thóc}
{Cần chuyển từ kho thứ hai sang kho thứ nhất $30$ tấn thóc để số thóc ở hai kho bằng nhau}
\loigiai{
\begin{itemchoice}
\item Đúng. Tổng số phần: $2 + 3 = 5\text{ (phần)}$.
\item Đúng. Kho 1: $150 : 5 \times 2 = 60\text{ (tấn)}$.
\item Đúng. Kho 2: $150 - 60 = 90\text{ (tấn)}$.
\item Sai. Hiệu số thóc là: $90 - 60 = 30\text{ (tấn)}$. Để hai kho bằng nhau chỉ cần chuyển $30 : 2 = 15\text{ (tấn)} 
eq 30\text{ tấn}$.
\end{itemchoice}
}
\end{ex}

\begin{ex} %[5D7V38]
Một mảnh vườn hình chữ nhật có chu vi $100\text{ m}$. Chiều dài bằng $\dfrac{3}{2}$ chiều rộng.
\choiceTF
{\True Nửa chu vi của mảnh vườn là $50\text{ m}$}
{\True Chiều dài của mảnh vườn là $30\text{ m}$}
{\True Chiều rộng của mảnh vườn là $20\text{ m}$}
{Diện tích của mảnh vườn là $500\text{ m}^2$}
\loigiai{
\begin{itemchoice}
\item Đúng. Nửa chu vi: $100 : 2 = 50\text{ (m)}$.
\item Đúng. Chiều dài: $50 : (3 + 2) \times 3 = 30\text{ (m)}$.
\item Đúng. Chiều rộng: $50 - 30 = 20\text{ (m)}$.
\item Sai. Diện tích: $30 \times 20 = 600\text{ (m}^2\text{)} 
eq 500\text{ m}^2$.
\end{itemchoice}
}
\end{ex}

\textbf{Phần 3. Trắc nghiệm trả lời ngắn}
\setcounterpageref{ex}{0}

\begin{ex} %[5D7V38]
Hai số có tổng bằng $120$. Tỉ số của hai số đó là $\dfrac{1}{3}$. Tìm số lớn.
\shortans{90}
\loigiai{
Tổng số phần bằng nhau là:
$1 + 3 = 4\text{ (phần)}$.\\
Số lớn là:
$120 : 4 \times 3 = 90$.
}
\end{ex}

\begin{ex} %[5D7V38]
Hiện nay tổng số tuổi của hai bố con là $48$ tuổi. Tuổi bố gấp $3$ lần tuổi con. Tính số tuổi của con hiện nay.
\shortans{12}
\loigiai{
Tổng số phần bằng nhau là:
$3 + 1 = 4\text{ (phần)}$.\\
Tuổi con hiện nay là:
$48 : 4 = 12\text{ (tuổi)}$.
}
\end{ex}

\begin{ex} %[5D7V38]
Hai thùng chứa tất cả $90\text{ l}$ dầu. Nếu chuyển $5\text{ l}$ dầu từ thùng thứ nhất sang thùng thứ hai thì số dầu ở thùng thứ nhất gấp đôi thùng thứ hai. Hỏi lúc đầu thùng thứ nhất chứa bao nhiêu lít dầu?
\shortans{65}
\loigiai{
Sau khi chuyển, tổng số dầu hai thùng vẫn là $90\text{ l}$.\\
Số dầu ở thùng thứ nhất sau khi chuyển là:
$90 : (2 + 1) \times 2 = 60\text{ (l)}$.\\
Lúc đầu thùng thứ nhất chứa:
$60 + 5 = 65\text{ (l)}$.
}
\end{ex}

\begin{ex} %[5D7V38]
Một lớp học có $36$ học sinh. Số học sinh nam bằng $\dfrac{4}{5}$ số học sinh nữ. Hỏi lớp học đó có bao nhiêu học sinh nữ?
\shortans{20}
\loigiai{
Tổng số phần bằng nhau là:
$4 + 5 = 9\text{ (phần)}$.\\
Số học sinh nữ của lớp là:
$36 : 9 \times 5 = 20\text{ (học sinh)}$.
}
\end{ex}

\begin{ex} %[5D7V38]
Chu vi một mảnh đất hình chữ nhật là $140\text{ m}$. Chiều rộng bằng $\dfrac{3}{4}$ chiều dài. Diện tích của mảnh đất đó là bao nhiêu mét vuông?
\shortans{1200}
\loigiai{
Nửa chu vi mảnh đất là:
$140 : 2 = 70\text{ (m)}$.\\
Chiều rộng là:
$70 : (3 + 4) \times 3 = 30\text{ (m)}$.\\
Chiều dài là:
$70 - 30 = 40\text{ (m)}$.\\
Diện tích mảnh đất là:
$40 \times 30 = 1\,200\text{ (m}^2\text{)}$.
}
\end{ex}

\begin{ex} %[5D7V38]
Ba bạn An, Bình, Cường có tất cả $120$ viên bi. Số bi của An bằng $\dfrac{1}{4}$ tổng số bi của cả ba bạn. Số bi của Bình bằng $\dfrac{1}{2}$ số bi của Cường. Hỏi Bình có bao nhiêu viên bi?
\shortans{30}
\loigiai{
Số bi của An là:
$120 \times \dfrac{1}{4} = 30\text{ (viên)}$.\\
Tổng số bi của Bình và Cường là:
$120 - 30 = 90\text{ (viên)}$.\\
Số bi của Bình là:
$90 : (1 + 2) \times 1 = 30\text{ (viên)}$.
}
\end{ex}
